\begin{statementbox}
	\textbf{\textcolor{CStatement}{条件}}
	\begin{description}[style=nextline, leftmargin=2.8em, labelsep=0.5em, itemsep=0.35em]
		\item[\textbf{\textcolor{CStatement}{条件 1（$n$ 重伯努利试验）：}}]
			随机变量 $X$ 表示 $n$ 次独立重复试验中成功的次数，每次试验成功概率为 $p$（$0 \le p \le 1$）。
	\end{description}
	\textbf{\textcolor{CStatement}{结论}}
	\begin{description}[style=nextline, leftmargin=2.8em, labelsep=0.5em, itemsep=0.45em]
		\item[\textbf{\textcolor{CStatement}{结论一（期望与方差）：}}]
			\textbf{\textcolor{CStatement}{直观描述：}}
			成功次数的均值与方差由 $n$ 和 $p$ 直接决定。
			\[
				E(X)=np,
				\qquad
				D(X)=np(1-p).
			\]
		\item[\textbf{\textcolor{CStatement}{推广结论（样本比例）：}}]
			\textbf{\textcolor{CStatement}{直观描述：}}
			成功频率的期望等于概率 $p$，方差随样本量增大而减小。
			\[
				E\!\left(\frac{X}{n}\right)=p,
				\qquad
				D\!\left(\frac{X}{n}\right)=\frac{p(1-p)}{n}.
			\]
	\end{description}
	\textbf{\textcolor{CStatement}{使用前提：}}
	各次试验相互独立，每次成功概率 $p$ 保持不变。
\end{statementbox}
